\documentclass[
% -- opções da classe memoir --
12pt,				% tamanho da fonte
openany,			% capítulos iniciam na próxima pág disponível (sem inserir páginas em branco)
oneside,			% para impressão em recto e verso. Oposto a oneside
a4paper,			% tamanho do papel. 
% -- opções da classe abntex2 --
%chapter=TITLE,		% títulos de capítulos convertidos em letras maiúsculas
%section=TITLE,		% títulos de seções convertidos em letras maiúsculas
%subsection=TITLE,	% títulos de subseções convertidos em letras maiúsculas
%subsubsection=TITLE,% títulos de subsubseções convertidos em letras maiúsculas
% -- opções do pacote babel --
english,			% idioma adicional para hifenização
french,				% idioma adicional para hifenização
spanish,			% idioma adicional para hifenização
brazil				% o último idioma é o principal do documento
]{abntex2}

% ---
% Pacotes básicos 
% ---
\usepackage{lmodern}			% Usa a fonte Latin Modern			
\usepackage[T1]{fontenc}		% Selecao de codigos de fonte.
\usepackage[utf8]{inputenc}		% Codificacao do documento (conversão automática dos acentos)
\usepackage{indentfirst}		% Indenta o primeiro parágrafo de cada seção.
\usepackage{color}				% Controle das cores
\usepackage{graphicx}			% Inclusão de gráficos
\usepackage{microtype} 			% para melhorias de justificação
% ---

% ---
% Pacotes adicionais, usados apenas no âmbito do Modelo Canônico do abnteX2
% ---
\usepackage{lipsum}				% para geração de dummy text
% ---
\usepackage{amssymb}   %para checkmark
% ---
% Pacotes de citações
% ---
\usepackage[brazilian,hyperpageref]{backref}	 % Paginas com as citações na bibl
\usepackage[alf]{abntex2cite}	% Citações padrão ABNT
% ---
% Pacotes adicionados pelos autores
% ---
\usepackage{enumitem}		% para personalizar listas numeradas

% --- 
% CONFIGURAÇÕES DE PACOTES
% ---

% Configurações do pacote backref
% Usado sem a opção hyperpageref de backref
\renewcommand{\backrefpagesname}{Citado na(s) página(s):~}
% Texto padrão antes do número das páginas
\renewcommand{\backref}{}
% Define os textos da citação
\renewcommand*{\backrefalt}[4]{
	\ifcase #1 %
	Nenhuma citação no texto.%
	\or
	Citado na página #2.%
	\else
	Citado #1 vezes nas páginas #2.%
	\fi}%

% Configuração do pacote enumitem
% Aplicando o recuo em listas numeradas globalmente
\setlist[enumerate]{leftmargin=2cm}
% ---

% ---
% Informações de dados para CAPA e FOLHA DE ROSTO
% ---
\titulo{Projeto Integrado de Extensão I}
\autor{BS Beauty Academy}
\local{São Paulo, SP - Brasil}
\data{2025}
\orientador{Marcelo Tavares}
\instituicao{%
	Instituto Federal de Educação, Ciência e\\ Tecnologia de São Paulo -- IFSP
	\par
	Tecnologia em Análise e Desenvolvimento de Sistemas
	\par
	Programa de graduação}
\tipotrabalho{Projeto Integrado de Extensão}
% O preambulo deve conter o tipo do trabalho, o objetivo, 
% o nome da instituição e a área de concentração 
\preambulo{Este projeto integrado de extensão, desenvolvido como parte do curso de Tecnologia em Análise e Desenvolvimento de Sistemas do Instituto Federal de Educação, Ciência e Tecnologia de São Paulo, tem como objetivo desenvolver uma aplicação web para otimizar a gestão e o agendamento de serviços de estética em um ambiente coworking. A área de concentração do projeto é a inovação tecnológica aplicada ao setor de serviços de beleza.}
% ---

% Configurações de aparência do PDF final

% alterando o aspecto da cor azul
\definecolor{blue}{RGB}{41,5,195}

% informações do PDF
\makeatletter
\hypersetup{
	%pagebackref=true,
	pdftitle={\@title}, 
	pdfauthor={\@author},
	pdfsubject={\imprimirpreambulo},
	pdfcreator={LaTeX with abnTeX2},
	pdfkeywords={abnt}{latex}{abntex}{abntex2}{trabalho acadêmico}, 
	colorlinks=true,       		% false: boxed links; true: colored links
	linkcolor=blue,          	% color of internal links
	citecolor=blue,        		% color of links to bibliography
	filecolor=magenta,      		% color of file links
	urlcolor=blue,
	bookmarksdepth=4
}
\makeatother
% --- 

% ---
% Posiciona figuras e tabelas no topo da página quando adicionadas sozinhas
% em um página em branco. Ver https://github.com/abntex/abntex2/issues/170
\makeatletter
\setlength{\@fptop}{5pt} % Set distance from top of page to first float
\makeatother
% ---

% ---
% Possibilita criação de Quadros e Lista de quadros.
% Ver https://github.com/abntex/abntex2/issues/176
%
\newcommand{\quadroname}{Quadro}
\newcommand{\listofquadrosname}{Lista de quadros}

\newfloat[chapter]{quadro}{loq}{\quadroname}
\newlistof{listofquadros}{loq}{\listofquadrosname}
\newlistentry{quadro}{loq}{0}

% configurações para atender às regras da ABNT
\setfloatadjustment{quadro}{\centering}
\counterwithout{quadro}{chapter}
\renewcommand{\cftquadroname}{\quadroname\space} 
\renewcommand*{\cftquadroaftersnum}{\hfill--\hfill}

\setfloatlocations{quadro}{hbtp} % Ver https://github.com/abntex/abntex2/issues/176
% ---

% --- 
% Espaçamentos entre linhas e parágrafos 
% --- 

% O tamanho do parágrafo é dado por:
\setlength{\parindent}{1.3cm}

% Controle do espaçamento entre um parágrafo e outro:
\setlength{\parskip}{0.2cm}  % tente também \onelineskip

% ---
% compila o indice
% ---
\makeindex
% ---

% ----
% Início do documento
% ----
\begin{document}
	
	% Seleciona o idioma do documento (conforme pacotes do babel)
	%\selectlanguage{english}
	\selectlanguage{brazil}
	
	% Retira espaço extra obsoleto entre as frases.
	\frenchspacing 
	
	% ----------------------------------------------------------
	% ELEMENTOS PRÉ-TEXTUAIS
	% ----------------------------------------------------------
	% \pretextual
	
	% ---
	% Capa
	% ---
	\imprimircapa
	% ---
	
	% ---
	% Folha de rosto
	% (o * indica que haverá a ficha bibliográfica)
	% ---
	\imprimirfolhaderosto*
	% ---
	
	% ---
	% Inserir a ficha bibliografica
	% ---
	
	% Isto é um exemplo de Ficha Catalográfica, ou ``Dados internacionais de
	% catalogação-na-publicação''. Você pode utilizar este modelo como referência. 
	% Porém, provavelmente a biblioteca da sua universidade lhe fornecerá um PDF
	% com a ficha catalográfica definitiva após a defesa do trabalho. Quando estiver
	% com o documento, salve-o como PDF no diretório do seu projeto e substitua todo
	% o conteúdo de implementação deste arquivo pelo comando abaixo:
	%
	% \begin{fichacatalografica}
		%     \includepdf{fig_ficha_catalografica.pdf}
		% \end{fichacatalografica}
	
	\begin{fichacatalografica}
		\sffamily
		\vspace*{\fill}					% Posição vertical
		\begin{center}					% Minipage Centralizado
			\fbox{\begin{minipage}[c][8cm]{13.5cm}		% Largura
					\small
					\imprimirautor
					%Sobrenome, Nome do autor
					
					\hspace{0.5cm} \imprimirtitulo  / \imprimirautor. --
					\imprimirlocal, \imprimirdata
					
					\hspace{0.5cm} \thelastpage p. : il. color; 30cm 
					
					\hspace{0.5cm} \imprimirorientadorRotulo~\imprimirorientador\\
					
					\hspace{0.5cm}
					\parbox[t]{\textwidth}{\imprimirtipotrabalho~--~\imprimirinstituicao,
						\imprimirdata.}\\
					
					\hspace{0.5cm}
					1. Graduação
					2. Extensão
					3. Integrado \\
					I. Marcelo Tavares.
					II. IFSP.
					III. Análise e Desenvolvimento de Sistemas. 
					IV. PIE1 			
			\end{minipage}}
		\end{center}
	\end{fichacatalografica}
	% ---
	
	% ---
	% Inserir errata
	% ---
	\begin{errata}
		pagina opcional para descrever algum correção que deve ser feita no documento.
		
		\begin{table}[htb]
			\center
			\footnotesize
			\begin{tabular}{|p{1.4cm}|p{1cm}|p{3cm}|p{3cm}|}
				\hline
				\textbf{Folha} & \textbf{Linha}  & \textbf{Onde se lê}  & \textbf{Leia-se}  \\
				\hline
				1 & 10 & auto-conclavo & autoconclavo\\
				\hline
			\end{tabular}
		\end{table}
		
	\end{errata}
	% ---
	
	% ---
	% Inserir folha de aprovação
	% ---
	
	% Isto é um exemplo de Folha de aprovação, elemento obrigatório da NBR
	% 14724/2011 (seção 4.2.1.3). Você pode utilizar este modelo até a aprovação
	% do trabalho. Após isso, substitua todo o conteúdo deste arquivo por uma
	% imagem da página assinada pela banca com o comando abaixo:
	%
	% \begin{folhadeaprovacao}
		% \includepdf{folhadeaprovacao_final.pdf}
		% \end{folhadeaprovacao}
	%
	\begin{folhadeaprovacao}
		
		\begin{center}
			{\ABNTEXchapterfont\large\imprimirautor}
			
			\vspace*{\fill}\vspace*{\fill}
			\begin{center}
				\ABNTEXchapterfont\bfseries\Large\imprimirtitulo
			\end{center}
			\vspace*{\fill}
			
			\hspace{.45\textwidth}
			\begin{minipage}{.5\textwidth}
				\imprimirpreambulo
			\end{minipage}%
			\vspace*{\fill}
		\end{center}
		
		Trabalho aprovado. \imprimirlocal, xx de xxx de 2025:
		
		\assinatura{\textbf{\imprimirorientador} \\ Orientador} 
		\assinatura{\textbf{Professor} \\ Convidado 1}
		\assinatura{\textbf{Professor} \\ Convidado 2}
		%\assinatura{\textbf{Professor} \\ Convidado 3}
		%\assinatura{\textbf{Professor} \\ Convidado 4}
		
		\begin{center}
			\vspace*{0.5cm}
			{\large\imprimirlocal}
			\par
			{\large\imprimirdata}
			\vspace*{1cm}
		\end{center}
		
	\end{folhadeaprovacao}
	% ---
	
	% ---
	% Dedicatória
	% ---
	\begin{dedicatoria}
		\vspace*{\fill}
		\centering
		\noindent
		\textit{ Este trabalho é dedicado às crianças adultas que,\\
			quando pequenas, sonharam em se tornar cientistas.} \vspace*{\fill}
	\end{dedicatoria}
	% ---
	
	% ---
	% Agradecimentos
	% ---
	\begin{agradecimentos}
	escrever agradecimentos (um paragrafo para cada?)
		
	\end{agradecimentos}
	% ---
	
	% ---
	% Epígrafe
	% ---
	\begin{epigrafe}
		\vspace*{\fill}
		\begin{flushright}
			\textit{``Não vos amoldeis às estruturas deste mundo, \\
				mas transformai-vos pela renovação da mente, \\
				a fim de distinguir qual é a vontade de Deus: \\
				o que é bom, o que Lhe é agradável, o que é perfeito.\\
				(Bíblia Sagrada, Romanos 12, 2)}
		\end{flushright}
	\end{epigrafe}
	% ---
	
	% ---
	% RESUMOS
	% ---
	
	% resumo em português
	\setlength{\absparsep}{18pt} % ajusta o espaçamento dos parágrafos do resumo
	\begin{resumo}
		o resumo deve ressaltar o
		objetivo, o método, os resultados e as conclusões do documento. A ordem e a extensão
		destes itens dependem do tipo de resumo (informativo ou indicativo) e do
		tratamento que cada item recebe no documento original. O resumo deve ser
		precedido da referência do documento, com exceção do resumo inserido no
		próprio documento. As palavras-chave devem figurar logo abaixo do
		resumo, antecedidas da expressão Palavras-chave:, separadas entre si por
		ponto e finalizadas também por ponto.
		
		\textbf{Palavras-chave}: latex. abntex. editoração de texto.
	\end{resumo}
	
	% resumo em inglês
	\begin{resumo}[abstract]
		\begin{otherlanguage*}{english}
			This is the english abstract.
			
			\vspace{\onelineskip}
			
			\noindent 
			\textbf{Keywords}: latex. abntex. text editoration.
		\end{otherlanguage*}
	\end{resumo}
	
	% ---
	% inserir lista de ilustrações
	% ---
	\pdfbookmark[0]{\listfigurename}{lof}
	\listoffigures*
	\cleardoublepage
	% ---
	
	% ---
	% inserir lista de quadros
	% ---
	\pdfbookmark[0]{\listofquadrosname}{loq}
	\listofquadros*
	\cleardoublepage
	% ---
	
	% ---
	% inserir lista de tabelas
	% ---
	\pdfbookmark[0]{\listtablename}{lot}
	\listoftables*
	\cleardoublepage
	% ---
	
	% ---
	% inserir lista de abreviaturas e siglas
	% ---
	\begin{siglas}
		\item [PIE1] Projeto Integrado de Extensão I
		\item [B2C] Business to Consumer
		\item [ECF] Emissor de Cupom Fiscal			\item [ERP] Enterprise Resource Planning (Sistema Integrado de Gestão Empresarial)
		\item [iOS] iPhone Operating System
		\item [MEI] Microempreendedor Individual
		\item [NFC- e] Nota Fiscal de Consumidor Eletrônica
		\item [PDV] Ponto de Venda
		\item [Pix] Pagamento Instantâneo
		\item [SAT] Sistema Autenticador e Transmissor de Cupons Fiscais
		\item [SLA] Service Level Agreement (Acordo de Nível de Serviço)
		\item [SMS] Short Message Service
		\item [TEF] Transferência Eletrônica de Fundos
		\item [TADS] Tecnologia em Análise e Desenvolvimento de Sistemas
	\end{siglas}
	% ---
	
	% ---
	% inserir lista de símbolos
	% ---
	\begin{simbolos}
			\item [R\$] Real (moeda brasileira)
			\item [\%]  Porcentagem
		
	\end{simbolos}
	% ---

	% inserir o sumario
	% ---
	\pdfbookmark[0]{\contentsname}{toc}
	\tableofcontents*
	\cleardoublepage

	% ----------------------------------------------------------
	\textual
	
% ----------------------------------------------------------
	\chapter{Introdução}
	
% o input do capitulo 1 já contem todas as seções dele	
%este arquivo contém todas as seções do capítulo de introdução

Segundo levantamento do Sebrae \cite{Sebrae_2025}
, em 2024 o setor de beleza no Brasil movimentou aproximadamente US\$\,27\,bilhões, colocando o país entre os cinco maiores mercados do mundo nesse ramo. Esse volume financeiro trouxe uma série de novidades e gerou, consequentemente, novas demandas. Diante de tantas mudanças e inovações, tornou-se indispensável que os empreendedores se adaptem rapidamente às tendências.

Conforme publicação do hub \emph{Beauty Fair} (Maior evento no setor da beleza no Brasil), até pouco tempo os profissionais autônomos precisavam deslocar-se até a residência de seus clientes para atendê-los ou firmar parcerias prestando serviços dentro de estabelecimentos de terceiros. Com o surgimento dos coworkings de beleza, esse cenário vem se transformando. A própria \emph{Beauty Fair} esclarece que um coworking de beleza é um espaço compartilhado que oferece infraestrutura para que profissionais da área possam trabalhar e colaborar. Trata-se de um local no qual cabeleireiros, maquiadores, esteticistas, manicures, massoterapeutas e demais especialistas podem alugar um posto de trabalho, dividir recursos e alcançar potenciais clientes \cite{BeautyFair}.

Reportagem online na Gazeta do Povo destaca que o ambiente coworking vem se consolidando como um dos modelos de negócio que mais crescem no Brasil, oferecendo ao profissional autônomo flexibilidade, troca de experiências e uma infraestrutura completa sem burocracia nem custos inesperados \cite{gazeta-coworking}. Nesse ambiente, o prestador de serviços tem o benefício de não precisar arcar com despesas de instalação ou manutenção de um salão próprio; basta utilizar o espaço, atender seus clientes e agendar a próxima sessão, preservando o controle sobre seus horários e ganhos.

À medida que esse formato de trabalho se expande, aumenta também a necessidade de maximizar a autonomia e a rentabilidade de cada profissional. Portanto, surge o desafio de gerir agendas, espaços e custos de forma ágil e intuitiva, evitando conflitos de reserva ou falhas de cobrança. Este projeto propõe-se a desenvolver uma aplicação web que atenda exatamente a essa demanda.

\section{Objetivos}

A aplicação web BS Beauty foi desenvolvida especialmente para gerenciar um salão de beleza que opera em modelo coworking, sob a gestão de nossa parceira de extensão Bruna. Seu objetivo principal é otimizar os processos internos e centralizar o agendamento de serviços, atendendo tanto às demandas da gestora quanto às necessidades dos profissionais autônomos.

\noindent\textbf{Para os Clientes Finais:} A plataforma possibilita o agendamento de serviços de forma intuitiva e flexível. Os clientes poderão escolher profissionais específicos ou optar pelo melhor horário disponível, visualizando facilmente a lista de prestadores, seus serviços, preços, tempo de execução e agendas atualizadas.

\noindent\textbf{Para os Profissionais Autônomos:} O sistema BS Beauty tem como propósito reforçar a autonomia dos profissionais sobre sua agenda e finanças. A aplicação permite bloquear horários, editar preços e a duração dos serviços, além de acompanhar os agendamentos realizados (sejam eles do dia, futuros ou passados) e visualizar relatórios detalhados com a receita gerada pelos serviços prestados.

\noindent\textbf{Para a Gestora:} Nossa parceira, Bruna, terá acesso a funcionalidades exclusivas que incluem análise de gastos (como luz e água), gerenciamento do aluguel ou comissão de cada profissional, visualização do fluxo de agendamentos em períodos específicos, envio de mensagens de marketing e promoções aos clientes, e acesso a relatórios financeiros detalhados. Ademais, a gestora poderá incluir ou remover profissionais da plataforma conforme a necessidade.
% ----------------------------------------------------------
\section{Problema e Solução Proposta}

RASCUNHO

A gestão de um local por empreendedores pequenos é frequentemente um desafio. Surgem demandas e processos que muitas vezes começam de froma manual, até que algo começa a dar errado e ve-se a necessidade de uma solução digital, a fim de diminuir erros lógicos e o esforço colocado na administração do negócio. No nosso caso, o objetivo geral da nossa aplicação é suprir as necessidades de um salão de beleza em modelo coworking que, como explicado anteriormete, é um modeol de trabalho recente (desde a pandemia) que visa atender as necessidades de diferentes autonomos, e não uma equipe com um objetivo em comum. O problema é como gerenciar a ocupação de cada profissional no local de trabalho além de controlar finanças e clientes. 

Nossa parceira, Bruna, já utilizava um sistema pra gerenciamento de seu salão, chamado Avec. Porém, apesar de ela já ter uma solução digitalizada, o sistema ainda deixava desejar em alguns aspectos (citar os pontos de insatifasção da bruna). 

Por isso, nossa solução é criar uma aplicação web que continue com todas as funcionalidades que já funcionam para nossa parceira Bruna (exemplificar), adicionar outras que fazem falta (exemplificar) e modificar requisitos funcionais ou não funcionais cuja ideia é boa, mas não funciona direito (como login).
% ----------------------------------------------------------
\section{Justificativa}
graficos com numeros, expor a relevância da solução - extensao e importancia
% ----------------------------------------------------------
\section{Análise da Concorrência}

Foi conduzida uma pesquisa de mercado centrada em plataformas brasileiras que combinam agendamento on-line e gestão financeira para espaços de beleza no modelo coworking. Deste levantamento emergiram três empresas que servirão de referência nesta análise: uma já amplamente consolidada no mercado nacional — embora atue além do universo coworking — e outras duas que, apesar de conhecidas, ainda estão em expansão, mas com foco mais relacionado ao da nossa proposta, o que as torna concorrentes que merecem maior atenção estratégica.

\subsection{Trinks}

%inicio de figura
\begin{figure}[htb]
	\centering
	\includegraphics[width=0.5\textwidth]{cap01-Introducao/Images/1.4.1_Trinks}
	\caption{Logo plataforma Trinks}
	\label{fig:Trinks}
\end{figure}

Trinks é uma plataforma já bem consolidada no mercado de gestão de negócios de beleza, com soluções
personalizadas para barbearias, salões de beleza e clínicas de estética. Criada em 2012, é hoje a
plataforma de gestão para beleza com a maior base instalada do país, englobando aproximadamente
2,8\,milhões de usuários e mais de 40\,mil estabelecimentos, sediada no Rio de Janeiro. A plataforma começou como um empreendimento de consultoria em software personalizado, mas logo identificou uma oportunidade no mercado da beleza e mudou de nicho. Em 2024, foi adquirida pelo grupo Stone, o que alavancou ainda mais funcionalidades do aplicativo, como o autoatendimento.
Atualmente, a Trinks oferece software de back-office (conjunto de módulos internos que controlam o funcionamento do negócio como finanças, estoque, comissões e relatórios), marketplace B2C e meios de pagamento próprios (Trinks Pay), funcionando praticamente como um “ERP + iFood” para salões e barbearias. Existe um
plano grátis que engloba apenas 150 agendamentos por mês, e os planos pagos variam de R\$ 59 a R\$ 249/mês \cite{Trinks}.

Além dos serviços comuns, seus principais diferenciais são:

\begin{itemize}
	\item Ponto de venda (PDV) completo: integração TEF, Pix e split de comissão, atendendo desde
	MEIs até redes com exigência de NFC-e e SAT/ECF;
	\item Marketplace \textit{Trinks.com}, que gera maior fluxo de clientes, expõe o salão ao
	público final e permite pagamento antecipado;
	\item Estrutura em nuvem madura, com SLA de 99{,}9\,\% e aplicativos nativos para
	iOS/Android.
\end{itemize}

Apesar dos grandes benefícios, identificamos algumas brechas do ponto de vista do negócio da nossa
parceira de extensão, Bruna:

\begin{itemize}
	\item A interface pode ser considerada “poluída” para clientes iniciantes, devido ao grande
	número de funcionalidades;
	\item Há pouco foco no aluguel de estações típico do coworking, exigindo ajustes manuais de
	comissão;
	\item Maior parte das funcionalidades estão presentes apenas nos planos superiores.
\end{itemize}

\subsection{Gendo}

%inicio de figura
\begin{figure}[htb]
	\centering
	\includegraphics[width=0.4\textwidth]{cap01-Introducao/Images/1.4.2_Gendo}
	\caption{Logo plataforma Gendo}
	\label{fig:Gendo}
\end{figure}

Lançado em 2017 e sediado em Curitiba-PR, o Gendo se posiciona como um \emph{hub}\footnote{Hub: plataforma centralizada que integra agenda, PDV, finanças e pagamentos em um único ambiente, funcionando como “nó” que organiza os fluxos de dados do negócio.} de gestão 100\,\% em nuvem para negócios além do setor da beleza, como estética, saúde, bem-estar, pet-shop e mais recentemente, espaços em formato coworking. 
Atualmente mantém mais de 10 mil assinantes, com maior penetração nas regiões Sul e Sudeste do Brasil. Foi criado no modelo SaaS com o intuito de oferecer prontamente agenda on-line, automação de lembretes (e-mail/WhatsApp), módulo financeiro completo e integrações com gateways de pagamento (Stone, Cielo e Mercado Pago). Atualmente, os planos são somente pagos e variam de R\$ 32 a R\$ 293/mês, após 14 dias de teste gratuito \cite{Gendo}.

Seus principais diferenciais são:
\begin{itemize}
	\item Caixa do profissional: Módulo pensado para coworking, possibilitando débito automático de aluguel de estação e visualização dos ganhos de cada profissional;
	\item Aplicativo Gendo Pro (iOS e Android): permite ao profissional ver a agenda, acompanhar comissões, pedir saques e registrar fotos de antes e depois dos serviços;
	\item Relatórios instantâneos: exibem ticket médio, previsão de faturamento e dados de cancelamentos, com opção de exportar para Excel.
\end{itemize}


Já os maiores pontos de melhoria identificados são:
\begin{itemize}
	\item Base instalada ainda pequena, limitando o efeito de rede junto a grandes franquias;
	\item Dependência de gateways externos, o que adiciona custo extra ao \emph{split} \footnote{Split é a divisão automática do pagamento entre salão e profissional que, se feita por um gateway externo, gera uma taxa extra.};
	\item Relatórios fiscais avançados disponíveis apenas no plano Premium.
\end{itemize}

\subsection{Avec}

\begin{figure}[htb]
	\centering
	\includegraphics[width=0.4\textwidth]{cap01-Introducao/Images/1.4.3_Avec}
	\caption{Logo plataforma Avec}
	\label{fig:Avec}
\end{figure}

A Avec é, hoje, a principal concorrente do nosso projeto, pois a entidade parceira que motivou este trabalho utiliza essa plataforma para gerenciar seu salão de beleza em modelo coworking. Por esse motivo, ela foi adotada como referência: buscamos manter as funcionalidades que já funcionam bem na Avec e, ao mesmo tempo, acrescentar ou aprimorar recursos que ainda fazem falta para a nossa parceira.

Lançada em 2014 e sediada em São Paulo-SP, a Avec se apresenta como solução``360º'' para salões, barbearias, esmaltarias, spas e estúdios de tatuagem. A plataforma integra software de gestão, um sistema próprio de pagamentos (\emph{Avec Pay}) e um marketplace B2C que encaminha novos clientes aos estabelecimentos. Segundo a empresa, mais de 40 mil negócios utilizam o serviço no
Brasil. Também desenvolvida no modelo SaaS, a ferramenta oferece agenda on\-/line multiprofissional com confirmações via WhatsApp ou SMS, ponto de venda completo com TEF, Pix e \emph{split} interno de comissões, além de módulo financeiro integrado. Dispõe ainda de uma carteira digital empregada em pacotes pré\=/pagos, gift-cards
e cashback, e de dois aplicativos: o \emph{Avec}, voltado ao cliente final, e o \emph{Avec Pro}, destinado aos profissionais. Há um plano gratuito ``\emph{Avec Go}'' que inclui funções básicas e cobra apenas a taxa transacional, enquanto os planos pagos variam de R\$ 77 a
R\$249 por mês \cite{Avec}.


Com base no feedback da nossa entidade parceira, destacam-se quatro funcionalidades que
a plataforma \emph{Avec} executa bem:
\begin{itemize}
	\item \emph{Split} instantâneo de comissões, dispensando gateways externos;
	\item Marketplace B2C e aplicativo do cliente, que ampliam a visibilidade do salão e aumentam os agendamentos on-line;
	\item App \emph{Avec Pro} (iOS/Android), no qual o profissional acompanha agenda,
	comissões, saques e registra fotos de “antes e depois” dos serviços.
\end{itemize}

As principais brechas identificadas são:
\begin{itemize}
	\item módulos fiscais avançados (NF-e e SAT) disponíveis apenas nos planos superiores;
	\item dependência do hardware e das tarifas do próprio \emph{Avec Pay} para uso pleno do sistema;
	\item custos adicionais para envios em massa de SMS/WhatsApp em campanhas de marketing;
	\item instabilidade recorrente: o domínio eventualmente fica fora do ar.
\end{itemize}

%seção 1.4 quadro comparativo

\subsection{Quadro comparativo}

\begin{quadro}[htb]
	\caption{\label{frame:comparativo_concorrência}Comparação entre as plataformas concorrentes e a aplicação proposta}
	\footnotesize
	\setlength{\tabcolsep}{4pt}
	\begin{tabular}{|p{6.8cm}|c|c|c|c|}
		\hline
		\textbf{Recurso}                                   & \textbf{Trinks} & \textbf{Gendo} & \textbf{Avec} & \textbf{BS Beauty}\\ \hline 
		Aplicação \textit{web}  & — & \checkmark & \checkmark & \checkmark \\ \hline
		Foco no coworking  & — & \checkmark & — & \checkmark \\ \hline
		Controle de acesso para gestão  & — & \checkmark  & \checkmark & \checkmark \\ \hline
		Agendamento de serviços 100\% on-line & \checkmark & \checkmark & \checkmark & \checkmark \\ \hline
		Controle de conflitos de agenda                    & \checkmark & \checkmark & \checkmark & \checkmark \\ \hline
		Avaliação pós-serviço  & \checkmark & \checkmark  & \checkmark  & \checkmark \\ \hline
		Plataforma do \textit{cliente}                               & \checkmark & \checkmark & \checkmark & \checkmark \\ \hline
		Plataforma do \textit{profissional}                                        & \checkmark & \checkmark & \checkmark & \checkmark \\ \hline
		Confirmação automática (WhatsApp / SMS / e-mail)                    & \checkmark & \checkmark & \checkmark & \checkmark \\ \hline
		Cálculo de \emph{Split} de comissão      & \checkmark & \checkmark & \checkmark & — \\ \hline
		Pagamento on-line                 & \checkmark & — & \checkmark & \ — \\ \hline
		Marketplace B2C                         & \checkmark & — & \checkmark & \ — \\ \hline
		Marketing integrado (envio em massa de SMS/WhatsApp/e-mail)  & \checkmark & — & \checkmark & \checkmark \\ \hline
		Gift-card                     & \checkmark & — & \checkmark & \checkmark \\ \hline
		Programa de indicação  & \checkmark & — & \checkmark & \checkmark \\ \hline
		Relatório financeiro em tempo real    & \checkmark & \checkmark & \checkmark & \checkmark \\ \hline
		Login simplificado com integração Google            & \checkmark & \checkmark &  — & \checkmark \\ \hline
		Fotos “antes e depois” anexadas ao histórico do cliente             & — & \checkmark & \checkmark & \checkmark \\ \hline
		Plano gratuito disponível                                           & \checkmark & — & \checkmark & \checkmark \\ \hline
		Lista de aniversariantes para promoções  & \checkmark & \checkmark & — & \checkmark \\ \hline
		
	\end{tabular}
	\fonte{Produzido pelos autores}
\end{quadro}


	% --------------------------------------
	
	\chapter{Revisão da Literatura}
	explicar os topicos principais do projeto
	
	%o input do capitulo 2 já contem todas as seções dele
	
	%este arquivo contém todas as seções do capitulo 2

\section{Histórico do assunto}
%----------------------------------------------
\section{Atualidade do assunto}
%----------------------------------------------
\section{Outros contextos do assunto (opcional)}	
	%  ----------------------------------------------------------
	
	\chapter{GESTÃO DO PROJETO}

Neste capítulo de Gestão do Projeto são abordados tópicos como a organização da equipe (definindo as responsabilidades, papéis e atividades de cada integrante do grupo), a metodologia adotada para a gestão e desenvolvimento do projeto, bem como o repositório da aplicação.

\section{Organização da Equipe}

A equipe do presente projeto é composta por seis docentes do curso de graduação Superior de Tecnologia em Análise e Desenvolvimento de Sistemas do \gls{ifsp}, a saber:

\begin{itemize}
	\item Alyson César Fumagalli dos Santos Júnior
	\item Bruno de Almeida Fischer
	\item Eliel da Silva
	\item Giovanna Camille Silva Carvalho
	\item Henrique Santiago Pires
	\item Henrriky Jhonny de Oliveira Bastos
\end{itemize}

O grupo de trabalho foi formado logo no início da primeira disciplina de \gls{pie} por estudantes que já haviam realizado diversos outros trabalhos em conjunto e, portanto, estavam acostumados a trabalhar em equipe.

Visando uma transmissão de informações clara e centralizada, utilizou-se o \emph{Discord} \cite{discord-2025} como ferramenta de comunicação para realizar reuniões assíncronas; e o \emph{Notion} \cite{notion-2025} para organizar documentos, atribuir tarefas e monitorar o andamento do projeto (juntamente com o \emph{ProjectLibre} \cite{projectlibre-2025}).

\subsection{Responsabilidades / Papéis / Atividades}
\label{subsec:papeis-equipe}

Para cada integrante da equipe foi definido um papel contendo responsabilidades e atividades definidas com base em suas respectivas proficiências em diferentes áreas, a fim de distribuir as tarefas do projeto de maneira eficiente.

O Quadro \ref{frame:distribuicao-papeis} descreve a distribuição dos membros do grupo com base em seus respectivos papéis de uma forma mais geral.

\begin{quadro}[ht]
	\setlength{\tabcolsep}{3pt}
	\begin{center}
		\caption{\label{frame:distribuicao-papeis}Membros e seus respectivos papéis}
		\begin{tabular}{|l|c|c|m{3cm}|c|}
			\hline
			\textbf{Membro} & \textbf{Gestor} & \textbf{\textit{Tech Lead}} & \centering\textbf{Desenvolvedor \textit{Full Stack}} & \textbf{Analista de Documentação} \\
			\hline
			Alyson & \checkmark &  & \centering\checkmark &  \\
			\hline
			Bruno &  &  & \centering\checkmark &  \\
			\hline
			Eliel &  &  & \centering\checkmark &  \\
			\hline
			Giovanna &  &  &  & \checkmark \\
			\hline
			Henrique &  &  & \centering\checkmark &  \\
			\hline
			Henrriky &  & \checkmark & \centering\checkmark &  \\
			\hline
		\end{tabular}
		\fonte{Produzido pelos autores}
	\end{center}
\end{quadro}

\indent Assim, constata-se que a equipe conta com 1 (um) gestor, responsável pelo gerenciamento de todo o projeto; 1 (um) \textit{tech lead}, encarregado de guiar a equipe de desenvolvimento; 5 (cinco) desenvolvedores \textit{full stack} (com 2 deles desempenhando outros papéis paralelos) incumbidos por desenvolver o sistema em todas as suas etapas e 1 (uma) analista de documentação para supervisionar as documentações do projeto. 

O Quadro \ref{quad:membros-atividades} apresenta as atividades desempenhadas pelos membros da equipe nas diversas áreas que contemplaram o desenvolvimento do projeto.

\begin{quadro}[h]
	\setlength{\tabcolsep}{2pt}
	\begin{center}
		\caption{\label{quad:membros-atividades}Membros e suas atividades}
		\begin{tabular}{|l|c|m{2.5cm}|m{2cm}|c|m{2cm}|c|}
			\hline
			\textbf{Membro} & \textbf{\emph{Notion}} & \centering\textbf{\textit{Front-End Back-End}} & \centering\textbf{Banco de Dados} & \textbf{Documentação} & \centering\textbf{Diário de Bordo} & \textbf{\emph{ProjectLibre}}\\
			\hline
			Alyson & \checkmark & \centering\checkmark & \centering\checkmark & \checkmark & \centering\checkmark & \checkmark \\
			\hline
			Bruno & \checkmark & \centering\checkmark & \centering\checkmark & \checkmark & & \checkmark\\
			\hline
			Eliel & \checkmark & \centering\checkmark & \centering\checkmark & \checkmark & & \checkmark\\
			\hline
			Giovanna & \checkmark &  & \centering\checkmark & \checkmark & \centering\checkmark & \checkmark\\
			\hline
			Henrique & \checkmark & \centering\checkmark & \centering\checkmark & \checkmark & & \\
			\hline
			Henrriky & \checkmark & \centering\checkmark & \centering\checkmark & \checkmark & & \\
			\hline
		\end{tabular}
		\fonte{Produzido pelos autores}
	\end{center}
\end{quadro}
\section{Metodologias de Gestão e Desenvolvimento}

Para o projeto, foi adotada a metodologia Scrum \cite{Scrum} tanto voltada para a gestão quanto para o desenvolvimento do sistema.

\subsection{Scrum}

O \textit{framework} Scrum foi escolhido por ter sido bastante estudado em disciplinas anteriores e também devido à equipe já ter uma certa familiaridade em trabalhar com ele. A metodologia foi adaptada de forma a incorporar alguns elementos do Kanban \cite{Kanban} a fim de ter um sistema visual para monitorar as atividades em andamento.

Tendo em mente as responsabilidades provenientes do Scrum, o Quadro \ref{frame:papeis-scrum} descreve como os membros do grupo foram distribuídos seguindo os papéis da metodologia.

\begin{quadro}[ht]
	\setlength{\tabcolsep}{3pt}
	\begin{center}
		\caption{\label{frame:papeis-scrum}Papéis dos integrantes com base no Scrum}
		\begin{tabular}{|l|c|c|c|}
			\hline
			\textbf{Membro} & \textbf{\textit{Product Owner}} & \textbf{\textit{Scrum Master}} & \textbf{Equipe de Desenvolvimento} \\
			\hline
			Alyson &  & \checkmark & \checkmark \\
			\hline
			Bruno &  &  & \checkmark \\
			\hline
			Eliel &  &  & \checkmark \\
			\hline
			Giovanna &  &  & \checkmark \\
			\hline
			Henrique &  &  & \checkmark \\
			\hline
			Henrriky & \checkmark &  & \checkmark \\
			\hline
		\end{tabular}
		\fonte{Produzido pelos autores}
	\end{center}
\end{quadro}

A plataforma Notion \cite{Notion} teve um papel significativo na implementação do Scrum: por meio dela, criou-se um espaço de trabalho para registrar o \textit{product backlog}, os \textit{sprints} e suas respectivas tarefas.

Com essa ferramenta, ainda foi possível atribuir informações --- como prioridade, \textit{status}, prazos e responsáveis --- para cada atividade, além de criar visualizações para simular as colunas e cartões do Kanban, criando uma estrutura organizada e informativa que facilitou o gerenciamento do projeto.

\subsubsection{Sprints}

Texto
	
	%-------------------------------------------
	
	\chapter{Desenvolvimento do projeto}
	% ---
	\section{Escopo do Projeto}
	\subsection{Regras de Negócio}
	\subsection{Requisitos Funcionais}
	\subsection{Requisitos Não Funcionais}
	% ---
	\section{Histórias de usuário}
	se aplicável para o scrum
	\subsection{Descrição das Histórias de Usuário}
	% ---
	\section{Arquitetura}
	\subsection{Definições da arquitetura}
	\subsection{Diagrama da arquitetura}
	% ---
	\section{Tecnologias}
	\subsection{Front-End}
	\subsection{Back-End}
	\subsection{Banco de dados}
	\subsection{Infraestrutura}
	\subsubsection{Amazon Web Services (AWS)}
	exemplo, colocar todos, abrindo demais itens
	%---
	\section{Testes e Manutenibilidade}
	\subsection{Plano de Testes}
	\subsection{Análise Estática}
	\subsection{Testes Funcionais}
	\subsection{Testes Unitários}
	\subsection{Testes de Componente}
	\subsection{Testes de Integração}
	\subsection{Testes end-to-end}
	\subsection{Testes não Funcionais}
	\subsection{Testes de Performance}
	\subsection{Testes de Carga}
	\subsection{Testes de Configuração}
	\subsection{Testes Automatizados}
	\subsection{Logs}
	\subsection{Code Convention}
	%---
	\section{Segurança, Privacidade e Legislação}
	\subsection{Critérios de segurança e privacidade}
	\subsection{Observância à Legislação}
	%---
	\section{Modelo de Banco de Dados}
	\subsection{Modelo Entidade Relacionamento - MER}
	\subsection{DIagrama Entidade Relacionamento - DER}
	tabelas
	\subsection{Dicionário de Dados}
	%---   
	\section{Cronograma}
	pensar no projeto todo, não só MVP
	\subsection{Análise da Duração do Projeto}
	considerar o gerenciamento ágil
	% ----------------------------------------------------------
	
	\chapter{Viabilidade Financeira}
	
	mesmo usando uma hospedagem gratis (AWS), precisamos pesquisar uma paga para colocar na tabela de custos
	
	\section{Custos}
	
	\section{Receitas}
	\section{Cenário realista}
	\section{Cenário Otimista}
	\section{Cenário Pessimista}
	% ---------------------------------------------------------------
	
	\chapter{Considerações Finais}
	\section{Dificuldades, escolhas e Descartes}
	\cite{ibge1993}
	% ----------------------------------------------------------
	
	ver documento "abntex2-modelo-include-comandos" para dicas e tutoriais de como fazer tabelas, graficos, quadros, inserir imagens e documentos externos etc aqui. 
	% ---
	
	%Página de referencias
	\renewcommand{\bibname}{Referências Bibliográficas} %renomeia o título da página de referencias
	\bibliography{referencias-bibliograficas}
	\bibliographystyle{abntex2-alf}
	\bibliographystyle{abntex2-num}
	
	
	% ----------------------------------------------------------
	% Glossário
	% ----------------------------------------------------------
	%
	% Consulte o manual da classe abntex2 para orientações sobre o glossário.
	%
	%\glossary
	
	% ----------------------------------------------------------
	% Apêndices
	% ----------------------------------------------------------
	
	% ---
	% Inicia os apêndices
	% ---
	\begin{apendicesenv}
		
		% Imprime uma página indicando o início dos apêndices
		\partapendices
		
		% ----------------------------------------------------------
		\chapter{exemplo 1}
		% ----------------------------------------------------------
		
		materiais desenvolvidos pelo próprio autor do trabalho
		
	\end{apendicesenv}
	
	% ----------------------------------------------------------
	% Anexos
	% Inicia os anexos
	
	\begin{anexosenv}
		
		% Imprime uma página indicando o início dos anexos
		\partanexos
		
		% ---
		\chapter{exemplo 1}
		matérias de outras fontes que não o próprio autor do trabalho
		
	\end{anexosenv}
	
	%---------------------------------------------------------------------
	% INDICE REMISSIVO
	%---------------------------------------------------------------------
	\phantompart
	\printindex
	%---------------------------------------------------------------------
	
\end{document}
